% Solutions of 09-version-control.
\section*{Chapter \ref{ch:b1:version-control} --- Version Control with Git}

\begin{solution}{exo:b1:version-control:1}
\omoutput[5-7]{code/bachelor-1/09-version-control/outputs/exo_status.txt}{The status.}
\code{books.csv} is modified in the \omterm{def:b1:version-control:staging}{staging area} (the Stone Bridge line, first column \code{M}) and modified again in
the \omterm{def:b1:version-control:staging}{working tree} since it was staged (the Winter Map line, second column \code{M}). \code{notes.txt} was never added:
untracked.
\end{solution}

\begin{solution}{exo:b1:version-control:2}
\code{sort-by-year}: 346e, 9fdd, 8da0, 3e49. \code{genres}: 9fdd, 8da0, 3e49 (the fast-forward did not create a \omterm{def:b1:version-control:commit}{commit}).
\code{main} after the \omterm{def:b1:version-control:merge}{merge}: all six, 58e3, 88f4, 346e, 9fdd, 8da0 and 3e49, since every one is an \omterm{def:b1:version-control:commit-graph}{ancestor} of 58e3.
\end{solution}

\begin{solution}{exo:b1:version-control:3}
(a) The \omterm{def:b1:version-control:merge}{merge base} of \code{main} and \code{fix} is $X$, the tip of \code{main}: fast-forward, \code{main} moves to the tip
of \code{fix}. (b) The \omterm{def:b1:version-control:merge}{merge base} is still $X$, but the tip of \code{main} is now a later \omterm{def:b1:version-control:commit}{commit}: a \omterm{def:b1:version-control:merge}{three-way merge}, with a
\omterm{def:b1:version-control:merge}{merge} \omterm{def:b1:version-control:commit}{commit}. (c) The \omterm{def:b1:version-control:merge}{merge base} is $X$, the tip of \code{main} being an \omterm{def:b1:version-control:commit-graph}{ancestor} of \code{fix}: \code{fix} already
contains every \omterm{def:b1:version-control:commit}{commit} of \code{main}, and Git answers ``Already up to date'' without changing anything.
\end{solution}

\begin{solution}{exo:b1:version-control:4}
\omoutput[4-9]{code/bachelor-1/09-version-control/outputs/exo_staged.txt}{The two diffs.}
\code{git add} copied the file with the year 1997 into the \omterm{def:b1:version-control:staging}{staging area}; the change to 2012 came after, in the \omterm{def:b1:version-control:staging}{working
tree} only. \code{git commit} records the \omterm{def:b1:version-control:staging}{staging area}: the year 1997 and Harbour Lights still dated 2011.
\end{solution}

\begin{solution}{exo:b1:version-control:5}
Keep a copy of the file with both changes, remove the new last line, \omterm{def:b1:version-control:commit}{commit} the year, then put the copy back and \omterm{def:b1:version-control:commit}{commit}
the book.
\omoutput[12-23]{code/bachelor-1/09-version-control/outputs/exo_split.txt}{Two \omterm{def:b1:version-control:commit}{commits} of one change each.}
With an editor, one would stage a part of the file with \code{git add -p}, which asks, for each changed block, whether
to stage it.
\end{solution}

\begin{solution}{exo:b1:version-control:6}
\omoutput{code/bachelor-1/09-version-control/outputs/exo_markers.txt}{The \omterm{def:b1:version-control:commit}{commit} Git accepted.}
Git accepts it: \code{-a} stages the file as it is, markers included, and staging is how a conflict is declared resolved.
The \omterm{def:b1:version-control:commit}{commit} records a file of seven lines, three of them markers, and two versions of the same book: any program that
reads the catalogue now fails or counts the book twice, on every \omterm{def:b1:version-control:branch}{branch} that \omterm{def:b1:version-control:merge}{merges} this one. \texttt{git diff -{}-{}check}
between the \omterm{def:b1:version-control:merge}{merge} and its first parent reports the markers, with an \omterm{def:b1:the-shell:exit-status}{exit status} of 2; a test of the file's format would
also catch them.
\end{solution}

\begin{solution}{exo:b1:version-control:7}
\omcode{code/bachelor-1/09-version-control/exercises/bisect_test.sh}{1}{12}{A test with three verdicts}
A program that does not parse stops before running the test, with a status of 1, which \code{bisect run} would read as
``bad'': \omterm{def:b1:version-control:bisection}{bisection} would then blame a \omterm{def:b1:version-control:commit}{commit} for a failure unrelated to the bug it is looking for, and conclude wrongly.
The status 125 says ``no verdict'', and \omterm{def:b1:version-control:bisection}{bisection} tests a neighbouring \omterm{def:b1:version-control:commit}{commit} instead. Statuses above 127 abort the
\omterm{def:b1:version-control:bisection}{bisection}, which is why the script maps any unexpected status to 125. The test \code{tests/test_history.py} runs the
script on the four cases.
\end{solution}

\begin{solution}{exo:b1:version-control:8}
The \omterm{def:b1:version-control:commit-graph}{ancestors} of $G$ are $G, D, F, C, E, B, A$; those of $H$ are $H, E, C, B, A$. The common \omterm{def:b1:version-control:commit-graph}{ancestors} are $E, C, B, A$, and
$C$, $B$, $A$ are \omterm{def:b1:version-control:commit-graph}{ancestors} of $E$: the \omterm{def:b1:version-control:merge}{merge base} is $E$, not $C$. For $D$ and $H$, the common \omterm{def:b1:version-control:commit-graph}{ancestors} are $C, B, A$,
and the \omterm{def:b1:version-control:merge}{merge base} is $C$. A \omterm{def:b1:version-control:merge}{merge} of $H$ into $G$ compares each file in $E$, in $G$ and in $H$: what $F$ and $D$ changed
since $E$ is kept from the side of $G$, what $H$ changed from the side of $H$.
\end{solution}

\begin{solution}{exo:b1:version-control:9}
\omoutput[3-14]{code/bachelor-1/09-version-control/outputs/exo_two_merges.txt}{The graph after two \omterm{def:b1:version-control:merge}{merges}.}
The first \omterm{def:b1:version-control:merge}{merge} \omterm{def:b1:version-control:commit}{commit} has parents M1 and A1; the second has parents the first \omterm{def:b1:version-control:merge}{merge} and B1. Every \omterm{def:b1:version-control:branch}{branch} was made from
8da0, which is the \omterm{def:b1:version-control:merge}{merge base} of both \omterm{def:b1:version-control:merge}{merges}.
\end{solution}

\begin{solution}{exo:b1:version-control:10}
Call the candidates $c_1, \dots, c_n$: the first bad \omterm{def:b1:version-control:commit}{commit} is one of them, since $c_0$ is good and $c_n$ bad. Let
$T(m)$ be the largest number of tests \omterm{def:b1:version-control:bisection}{bisection} needs when $m$ consecutive candidates remain, the last of them known bad.
$T(1) = 0$. For $m \geq 2$, testing the candidate of position $\lfloor m / 2 \rfloor$ leaves the first $\lfloor m / 2
\rfloor$ candidates if it is bad and the last $\lceil m / 2 \rceil$ if it is good: $T(m) \leq 1 + T(\lceil m / 2
\rceil)$. By induction, $T(m) \leq \lceil \log_2 m \rceil$, since for $m \geq 2$, $\lceil \log_2 \lceil m / 2 \rceil
\rceil = \lceil \log_2 m \rceil - 1$. Lower bound: a method that makes at most $t$ tests, each with two answers, can end
in at most $2^t$ different ways, and it must name $n$ different culprits for the $n$ possible histories: $2^t \geq n$,
so $t \geq \lceil \log_2 n \rceil$. The test \code{tests/test_history.py} checks the recurrence for $n < 300$.
\end{solution}

\begin{solution}{exo:b1:version-control:11}
Let the candidates be the \omterm{def:b1:version-control:commit-graph}{ancestors} of the known bad \omterm{def:b1:version-control:commit}{commit} $b$ that are not \omterm{def:b1:version-control:commit-graph}{ancestors} of a good \omterm{def:b1:version-control:commit}{commit}, $b$ included;
the test picks a candidate $x \neq b$. If $x$ is bad, the candidates become those that are \omterm{def:b1:version-control:commit-graph}{ancestors} of $x$, and $b$,
which is not an \omterm{def:b1:version-control:commit-graph}{ancestor} of $x$ (the graph has no cycle), leaves; if $x$ is good, $x$ leaves. The set loses at least one
\omterm{def:b1:version-control:commit}{commit} at each test, so \omterm{def:b1:version-control:bisection}{bisection} ends, with $b$ alone. In the example, the candidates are $C_1, \dots, C_k$ and $M$, and
each $C_i$ has no candidate among its \omterm{def:b1:version-control:commit-graph}{ancestors} but itself: a ``good'' verdict removes one \omterm{def:b1:version-control:commit}{commit} only. If the culprit
is $M$, every $C_i$ must be tested: $k$ tests in the worst case, against $\lceil \log_2 (k + 1) \rceil$ on a linear history
of as many candidates (3 instead of 7 for $k = 7$). \omterm{def:b1:version-control:merge}{Merges} of many short \omterm{def:b1:version-control:branch}{branches} make \omterm{def:b1:version-control:bisection}{bisection} slower.
\end{solution}

\begin{solution}{exo:b1:version-control:12}
\omcode{code/bachelor-1/09-version-control/listings/touch_count.sh}{1}{8}{\omterm{def:b1:version-control:commit}{Commits} per file}
\omoutput{code/bachelor-1/09-version-control/outputs/exo_touch.txt}{The counts on the generated history.}
\texttt{-{}-{}name-{}only} lists the files changed by each \omterm{def:b1:version-control:commit}{commit} and \texttt{-{}-{}format=} prints nothing else, leaving blank lines
that \code{grep} removes; \texttt{-{}-{}no-{}merges} skips the \omterm{def:b1:version-control:merge}{merge}, which would otherwise count again files brought from the
\omterm{def:b1:version-control:branch}{branch}. All four counts follow from the generator: \code{books.csv} is created by the first \omterm{def:b1:version-control:commit}{commit} and extended by 38 of
the next 39 main \omterm{def:b1:version-control:commit}{commits}, the other one being the bad \omterm{def:b1:version-control:commit}{commit}: 39; \code{covers.csv} 23, one per \omterm{def:b1:version-control:branch}{branch} \omterm{def:b1:version-control:commit}{commit};
\code{catalogue.py} 2, its creation and the bad \omterm{def:b1:version-control:commit}{commit}; \code{check.py} 1. The test \code{tests/test_history.py} checks
them from the generator's constants.
\end{solution}

\begin{solution}{pb:b1:version-control:1}
\textbf{1.} 64: 40 \omterm{def:b1:version-control:commit}{commits} on \code{main}, 23 on the \omterm{def:b1:version-control:branch}{branch}, and the \omterm{def:b1:version-control:merge}{merge}.

\textbf{2.} The \omterm{def:b1:version-control:merge}{merge}, 8424896. \texttt{|\textbackslash} under it opens a second column, its second parent's line; the
\code{|/} near the bottom joins that column to the \omterm{def:b1:version-control:commit}{commit} the \omterm{def:b1:version-control:branch}{branch} was made from.

\textbf{3.} The \omterm{def:b1:version-control:commit}{commits} reached from \code{HEAD} through first parents only: the \omterm{def:b1:version-control:merge}{merge} and the 40 \omterm{def:b1:version-control:commit}{commits} of \code{main}.

\textbf{4.} The tip of \code{covers} has as \omterm{def:b1:version-control:commit-graph}{ancestors} its 23 \omterm{def:b1:version-control:commit}{commits} and the first 20 \omterm{def:b1:version-control:commit}{commits} of \code{main}: 43.

\textbf{5.} No: it is the 31st \omterm{def:b1:version-control:commit}{commit} of \code{main}, made after the \omterm{def:b1:version-control:branch}{branch} left at the 20th. Yes for \code{HEAD}, which
is reached from it through the \omterm{def:b1:version-control:merge}{merge}'s first parent.

\textbf{6.} The graph would be the same, since it is made of parent links, not dates. The identifiers are the same at
each run anyway, since they are hashes of contents, names, dates and parents that the seed and the fixed dates
determine; with dates taken from the clock, they would change at every run.

\textbf{7.} 66ce86b, ``Add The Copper Map'', the 20th \omterm{def:b1:version-control:commit}{commit} of \code{main}: the test
\code{if i == FORK:} of the loop makes the \omterm{def:b1:version-control:branch}{branch} just after it.

\textbf{8.} The \omterm{def:b1:version-control:commit}{commits} reachable from \code{HEAD} and not from \code{covers}: the 20 \omterm{def:b1:version-control:commit}{commits} 21 to 40 of \code{main},
and the \omterm{def:b1:version-control:merge}{merge}.

\textbf{9.} The 23 \omterm{def:b1:version-control:commit}{commits} of the \omterm{def:b1:version-control:branch}{branch} are not \omterm{def:b1:version-control:commit-graph}{ancestors} of the first parent of the \omterm{def:b1:version-control:merge}{merge}: the \omterm{def:b1:version-control:merge}{merge} brought exactly
them into \code{main}.

\textbf{10.} \code{catalogue.py}: changed on \code{main} only (\code{min} to \code{max}), so the \omterm{def:b1:version-control:merge}{merge} keeps
\code{max}. \code{books.csv}: on \code{main} only, so \code{main}'s version. \code{covers.csv}: created on the \omterm{def:b1:version-control:branch}{branch}
only, so the \omterm{def:b1:version-control:branch}{branch}'s version.

\textbf{11.} No file was changed on both sides since the \omterm{def:b1:version-control:merge}{merge base}.

\textbf{12.} The tip of \code{covers} has the catalogue of \omterm{def:b1:version-control:commit}{commit} 20, with \code{min}: it passes. The \omterm{def:b1:version-control:merge}{merge} has
\code{max}: it fails.

\textbf{13.} A program can break through a change to another file (the data, a library, a configuration), and a file may
be changed by dozens of \omterm{def:b1:version-control:commit}{commits}, of which \code{log} cannot say which one broke the test.

\textbf{14.} \code{min} became \code{max}: \code{oldest} returns the book with the largest year, and among the three
books of the test that is Harbour Lights, 2011.

\textbf{15.} That \code{HEAD} is bad and the first \omterm{def:b1:version-control:commit}{commit} good. The candidates are the \omterm{def:b1:version-control:commit-graph}{ancestors} of \code{HEAD} that are
not \omterm{def:b1:version-control:commit-graph}{ancestors} of the first \omterm{def:b1:version-control:commit}{commit}: 63 \omterm{def:b1:version-control:commit}{commits}.

\textbf{16.} Bad, good, good, good, good, bad. Git announces 31, 15, 7, 3, 1 and 0 \omterm{def:b1:version-control:commit}{commits} left to test after each
choice: the candidates halve at every test.

\textbf{17.} The first test, at ``Add The Amber Map'', the 32nd \omterm{def:b1:version-control:commit}{commit} of \code{main}, was bad: the candidates became its
\omterm{def:b1:version-control:commit-graph}{ancestors}, among which no \omterm{def:b1:version-control:commit}{commit} of the \omterm{def:b1:version-control:branch}{branch}.

\textbf{18.} With a syntax error, Python stops with status 1 before running the test, and \code{bisect run} reads ``bad''.
Had such a \omterm{def:b1:version-control:commit}{commit} been tested, \omterm{def:b1:version-control:bisection}{bisection} would have treated it as broken by the bug: had it been on the \omterm{def:b1:version-control:branch}{branch}, before
the first test removed the \omterm{def:b1:version-control:branch}{branch}, it could have been named the first bad \omterm{def:b1:version-control:commit}{commit}. \code{bisect_test.sh} returns 125 for it,
and the \omterm{def:b1:version-control:commit}{commit} is skipped.

\textbf{19.} By induction: after the test $k + 1$, at most $\lceil m_k / 2 \rceil$ candidates remain, and
$\lceil \lceil n / 2^k \rceil / 2 \rceil = \lceil n / 2^{k + 1} \rceil$.

\textbf{20.} $\lceil 63 / 2^5 \rceil = 2$ and $\lceil 63 / 2^6 \rceil = 1$: six tests.

\textbf{21.} Along \code{main} from the oldest, the \omterm{def:b1:version-control:commit}{commits} 2 to 31: 30 tests, the last one bad. In the worst case, 62
tests for 63 candidates, since the last candidate need not be tested.

\textbf{22.} After choosing the first \omterm{def:b1:version-control:commit}{commit}, 31 candidates remain whatever its verdict, and about $\log_2 31 \approx 5$
tests are still needed after this one: five, plus the one about to run, make six.

\textbf{23.} The bad \omterm{def:b1:version-control:commit}{commit} lies on \code{main} after the fork, and the first test, bad, left only \omterm{def:b1:version-control:commit}{commits} of \code{main}:
from then on the candidates formed a linear history, on which each test halves them (\cref{exo:b1:version-control:10}).
Had the first verdict been good, the \omterm{def:b1:version-control:branch}{branch} would have stayed, with the risk shown in \cref{exo:b1:version-control:11}.

\textbf{24.} \emph{\omterm{def:b1:version-control:bisection}{Bisection} finds the culprit among 64 \omterm{def:b1:version-control:commit}{commits} in at most 6 tests.} With the first \omterm{def:b1:version-control:commit}{commit} good and the
last bad, the 63 candidates halve at each test, so $\lceil \log_2 63 \rceil = 6$ tests suffice; the captured run of
\code{git bisect run} took exactly six, and named ``Tidy the catalogue functions'', the \omterm{def:b1:version-control:commit}{commit} that changed \code{min}
into \code{max}, which the chapter's test confirms by running \code{check.py} on all 64 \omterm{def:b1:version-control:commit}{commits}.
\end{solution}
